% !TEX root = ../../main.tex
% ============================================================================
% Part IV; Platonic Introduction.
%
% Part IV has six typed carriers: the Vol I boundary-holographic datum,
% the associator face scheme, the celestial Mellin-shadow comparison,
% the Monster orbifold-Hochschild carrier, the K3/BKM arithmetic
% comparison surface, and the Vol II $\mathsf{SC}^{\mathrm{ch,top}}$
% to $E_3^{\mathrm{top}}$ topologisation map.  The source theorem is a
% statement about these carriers and their comparison maps, not an
% identification of all physical realisations with one object.
% ============================================================================

\providecommand{\cA}{\mathcal{A}}
\providecommand{\cB}{\mathcal{B}}
\providecommand{\cH}{\mathcal{H}}
\providecommand{\cM}{\mathcal{M}}
\providecommand{\cW}{\mathcal{W}}
\providecommand{\Walg}{\mathcal{W}}
\providecommand{\fg}{\mathfrak{g}}
\providecommand{\fgrt}{\mathfrak{grt}_1}
\providecommand{\fsl}{\mathfrak{sl}}
\providecommand{\Vir}{\mathrm{Vir}}
\providecommand{\GRT}{\mathrm{GRT}_1}
\providecommand{\Q}{\mathbb{Q}}
\providecommand{\Z}{\mathbb{Z}}
\providecommand{\ChirAlg}{\mathrm{ChirAlg}}
\providecommand{\Einfty}{E_\infty}
\providecommand{\Eone}{E_1}
\providecommand{\Etwo}{E_2}
\providecommand{\Ethree}{E_3}
\providecommand{\MC}{\mathrm{MC}}
\providecommand{\Face}{\mathrm{Face}}
\providecommand{\Phol}{\Phi_{\mathrm{hol}}}
\providecommand{\Zder}{\mathcal{Z}^{\mathrm{der}}_{\mathrm{ch}}}
\providecommand{\Holbd}{\mathsf{Hol}^{\mathrm{bd}}}
\providecommand{\SCchtop}{\mathsf{SC}^{\mathrm{ch,top}}}
\providecommand{\FactAlg}{\mathsf{FactAlg}}
\providecommand{\OC}{\mathrm{OC}}
\providecommand{\VOA}{\mathrm{VOA}}
\providecommand{\DW}{\mathrm{DW}}
\providecommand{\Res}{\operatorname{Res}}
\providecommand{\Mel}{\mathrm{Mel}}
\providecommand{\Sh}{\mathrm{Sh}}

\chapter{The Modular Tower's typed witnesses and physics bridges}
\label{ch:part-iv-platonic-introduction}

Part~\ref{part:modular-tower} develops Theorem~D
$\mathrm{obs}_g(A_b) = \kappa(A_b)\cdot\lambda_g$ on the
uniform-weight scalar lane, the shadow-tower quadrichotomy via
Borel--Riccati spectral analysis, and the modular trace + clutching
package on the open category. The present chapter inscribes the five
typed comparison carriers — boundary-holographic datum, associator
face scheme, celestial Mellin-shadow, Monster orbifold-Hochschild
carrier, K3/BKM arithmetic comparison surface, and Vol~II
$\mathsf{SC}^{\mathrm{ch,top}}$ bulk — that realise the
modular-tower obstruction class $\mathrm{obs}_g$ in concrete physics
computations (Feynman BV, Kontsevich weights, BRST, arithmetic
shadows).

Let $\cA$ be a Koszul-compatible $\Eone$-chiral algebra on a smooth
algebraic curve~$X$.  The Vol~I source is the boundary-holographic
datum
\begin{equation}
 \Phol_X(\cA) \;=\;
 \Bigl(\bar{B}^{\mathrm{ch}}(\cA),\Theta_\cA,r_\cA(z),
 \Zder(\cA),\nabla^{\mathrm{hol}}_\cA\Bigr).
 \label{eq:part-iv-holographic-functor}
\end{equation}
Here $\bar{B}^{\mathrm{ch}}(\cA)$ is the conilpotent bar coalgebra,
$\Theta_\cA$ is the universal Maurer--Cartan element,
$r_\cA(z)=\Res^{\mathrm{coll}}_{0,2}(\Theta_\cA)$ is the binary
collision residue, and
$\Zder(\cA)=C^\bullet_{\mathrm{ch}}(\cA,\cA)$ is the cochain-level
derived chiral centre.  The last object is the closed-sector algebra
internal to the boundary theory.  It is identified with a physical
bulk only after an open--closed comparison quasi-isomorphism
\[
 \OC_\cA:\mathcal O^{\mathrm{phys}}_{\mathrm{bulk}}(\cA)
 \xrightarrow{\;\sim\;} \Zder(\cA)
\]
has been constructed.

The five comparison carriers are:
\[
\begin{array}{c|c|c}
\text{carrier} & \text{map or object} & \text{mathematical content} \\
\hline
\text{GRT face scheme} &
\Face_K(\cA) &
\GRT_1(K)\text{-action on associator choices} \\
\text{celestial factorisation} &
\Sh\Mel\,\Phol_X(\cA^{\mathrm{soft}}_{\mathcal T})
\to \mathcal V^{\mathrm{cel}}_{\mathcal T} &
\text{soft-collinear comparison on }\CP^1 \\
\text{Monster orbifold} &
V_{\Lambda_{24}}^+\oplus V_{\Lambda_{24}}^{T,+}=V^\natural &
\Z/2\text{-crossed VOA extension and }H^3(\Z/2,U(1)) \\
\text{K3/BKM arithmetic} &
\mathbf H_{\Delta_5}\leadsto
\mathcal D_\hbar(\mathcal Y^{\mathrm{Hall}}_\hbar(\mathrm{CoHA}_{K3\times E})) &
\text{Vol III comparison surface} \\
\text{topologisation} &
\SCchtop(\cA,\Zder(\cA))\to E_3^{\mathrm{top}}(\Zder(\cA)) &
\text{Vol II conformal-vector/completion map.}
\end{array}
\]
The theorem surface of Part~\ref{part:physics-bridges} is the
compatibility of these carriers with the bar datum
\eqref{eq:part-iv-holographic-functor}.  None of the comparison
carriers is a substitute for the bar coalgebra, the Verdier dual
$\cA^!$, or the derived centre.

\section{The six carriers}
\label{sec:part-iv-deficiency}

The carrier separation is forced by type.  The bar coalgebra
$\bar{B}^{\mathrm{ch}}(\cA)=T^c(s^{-1}\bar\cA)$ classifies twisting
morphisms.  The Koszul dual algebra
$\cA^!=(H^\bullet\bar{B}^{\mathrm{ch}}(\cA))^\vee$ is obtained after
Verdier duality on the Koszul locus.  The derived chiral centre
$\Zder(\cA)$ is a Hochschild cochain object.  The physical bulk of an
open--closed field theory is a fourth object until the map
$\OC_\cA$ is supplied.

The GRT carrier acts on associator choices in the bar differential.
The celestial carrier applies the Mellin and shadow transforms to the
soft-sector collision residue.  The Monster carrier is the
$\Z/2$-crossed orbifold extension of the Leech lattice VOA.  The
K3/BKM carrier belongs to the Vol III Hall--Borcherds comparison
surface.  The topologisation carrier is the Vol II passage
\[
 \SCchtop(\cA,\Zder(\cA))
 \xrightarrow{\;T_{\mathrm{conf}},\,\widehat{\mathrm{wt}}\;}
 E_3^{\mathrm{top}}(\Zder(\cA)^\wedge),
\]
defined only after the conformal-vector and weight-completion
hypotheses used in Vol II are present.

\section{The boundary-holographic functor}
\label{sec:part-iv-holographic-functor}

The functor $\Phol_X$ assembles the bar coalgebra, the universal
Maurer--Cartan element, the classical r-matrix, the derived chiral
centre, and the Gauss--Manin connection into a boundary datum.

\begin{definition}[Boundary-holographic datum]
\label{def:part-iv-chiral-holographic-datum}
A \emph{boundary-holographic datum} over a smooth algebraic curve~$X$
is a tuple
$\mathscr{H} = (\cB, \Theta, r, Z, \nabla^{\mathrm{hol}})$
where:
\begin{itemize}[leftmargin=*]
 \item $\cB$ is a conilpotent $\Eone$-chiral coassociative coalgebra
 over~$X$ with coproduct $\Delta$ of cohomological degree~$0$ and
 differential $d_\cB$ of cohomological degree~$+1$;
 \item $\Theta \in \Hom^1(\cB, \cB)^{\mathrm{prim}}$ is a
 primitive Maurer--Cartan element:
 $d_\cB \Theta + \tfrac{1}{2} [\Theta, \Theta] = 0$;
 \item $r = \Res^{\mathrm{coll}}_{0,2}(\Theta) \in \cA^! \otimes
 \cA^! [\![z^{-1}]\!]$ is its degree-$2$ collision residue;
 \item $Z$ is a brace dg algebra, to be realised for a chiral algebra
 $\cA$ as the derived chiral centre
 $\Zder(\cA)=C^\bullet_{\mathrm{ch}}(\cA,\cA)$;
 \item $\nabla^{\mathrm{hol}}: \cB \to \cB \otimes \Omega^1_X$ is a
 flat Gauss--Manin-type connection compatible with~$\Theta$ and inducing
 the Gauss--Manin connection on $Z$ in families.
\end{itemize}
The category of boundary-holographic data over~$X$ is denoted
$\Holbd_X$.  A physical open--closed realisation is extra structure:
a quasi-isomorphism
\[
 \OC_\cA:\mathcal O^{\mathrm{phys}}_{\mathrm{bulk}}(\cA)
 \xrightarrow{\sim} Z
\]
from the physical bulk observables of the relevant field theory to
the derived chiral centre.  The datum $\mathscr H$ does not contain
this map by notation.
\end{definition}

\begin{definition}[The boundary-source functor]
\label{def:part-iv-Phol}
The \emph{boundary-source functor}
\[
 \Phol_X:\;\; \mathsf{ChirAlg}^{\mathrm{Kosz}}_X
 \;\longrightarrow\;
 \Holbd_X
\]
sends a Koszul-compatible $\Eone$-chiral algebra $\cA$ on~$X$ to the
tuple $\Phol_X(\cA) = (\bar{B}^{\mathrm{ch}}(\cA), \Theta_\cA,
r(z), \Zder(\cA), \nabla^{\mathrm{hol}}_\cA)$, where
$\Theta_\cA = D_\cA - d_0$ is the universal Maurer--Cartan
element \textup{(}Theorem~\ref{thm:mc2-bar-intrinsic}\textup{)},
$r(z)$ is its degree-$2$ collision residue,
$\Zder(\cA) = \mathrm{ChirHoch}^\bullet(\cA)$ is the chiral
Hochschild complex, and $\nabla^{\mathrm{hol}}_\cA$ is the Gauss--Manin
connection on the bar fibres over the configuration space.
\end{definition}

The projection onto any one coordinate recovers a canonical
construction: projection onto~$\cB$ alone is the
bar construction; projection onto $(\cB, \Theta)$ is the
Maurer--Cartan moduli MC\_2; projection onto~$r(z)$ is the classical
r-matrix hub; projection onto $\Zder(\cA)$ is the derived centre of
$E_n$ Koszul duality (Chapter~\ref{ch:en-koszul-duality}), not a
physical bulk before an open--closed map $\OC_\cA$ is given;
projection onto $\nabla^{\mathrm{hol}}_\cA$ at critical level is
the algebra of opers (Chapter~\ref{ch:derived-langlands}).

The coordinates of $\Phol_X$ assemble into a typed operadic passage
\[
 \SCchtop(\cA,\Zder(\cA))
 \longrightarrow
 E_3^{\mathrm{top}}(\Zder(\cA)^\wedge)
\]
only under the conformal-vector and completion hypotheses used in
the Vol II topologisation theorem.  The generic carrier is the
two-coloured object $\SCchtop(\cA,\Zder(\cA))$.
Each arrow of the passage is a distinct mathematical construction:
restriction to a codimension-two defect; ordered bar complex;
categorified averaging; higher Deligne. The passage is proved
single-direction, arrow by arrow, and should not be read as a
round-trip isomorphism. Closure to a full operadic circle requires a
conformal vector making translations $Q$-exact and is proved
chain-level on the original complex only for the affine
Kac--Moody case at non-critical level
\textup{(}by the Sugawara construction\textup{)};
for general chiral algebras with conformal vector, the closure is
realised at the cohomological level via the iterated Sugawara
theorem of Vol~II
\texttt{V2-thm:iterated-sugawara-construction}, and remains
conjectural at chain level.
Without such a chosen conformal vector, or at loci where the
topologisation comparison fails \textup{(}critical level and
nonconformal free-field presentations\textup{)}, the intermediate object is
$\SCchtop$ rather than $E_3^{\mathrm{top}}$:
the Swiss-cheese two-coloured chiral-topological operad is the
correct home, and Dunn additivity does not apply at the coloured
level. The Costello--Gwilliam factorization algebra $F^{\mathrm{CS}}_k$
on $\R \times \C$ gives a single-coloured factorization envelope of
the boundary chiral algebra, distinct from
$\mathsf{SC}^{\mathrm{ch,top}}$; the two presentations agree on the
affine Kac--Moody locus at non-critical level and diverge elsewhere.

\section{$\Face_K(\cA)$ as associator carrier}
\label{sec:part-iv-face-torsor}

The bar differential $d_\cB$ and the universal Maurer--Cartan
element $\Theta_\cA$ are defined in terms of a choice of Drinfeld
associator $\Phi \in \mathsf{Assoc}(K)$ over a coefficient field~$K$.
For rational associators one may take $K=\Q$; for the KZ/Brown
motivic face one must work over the corresponding motivic-period
realisation field. Different choices of
associator yield isomorphic but not equal holographic data; the
ambiguity is precisely the Drinfeld Grothendieck--Teichmuller
group~$\GRT_1(K)$ after a formality base point has been fixed. The
rigidified orbit of $\Phol_X(\cA)$ under the $\GRT_1(K)$-action is the
\emph{face scheme}
\begin{equation}
 \Face_K(\cA) \;=\; \bigl\{\,\Phol_X(\cA,\Phi)\;:\;
 \Phi \in \mathsf{Assoc}(K)\,\bigr\}\big/\mathrm{iso}.
 \label{eq:part-iv-face-scheme}
\end{equation}
The fingerprint map is a further quotient of this carrier. Its named
charts are comparison maps from the same collision residue
$r_\cA(z)$ to concrete presentations; the list of named charts is not
the full $\GRT_1(K)$-orbit.

\section{Celestial factorization comparison}
\label{sec:part-iv-celestial}

The celestial factorization algebra $\mathcal{V}^{\mathrm{cel}}$ on the
celestial sphere $\CP^1$ of four-dimensional asymptotic null infinity
arises from the leading soft theorems of a bulk four-dimensional
gauge theory: a Mellin transform of the S-matrix followed by a
shadow transform of the conformal-primary basis yields a
two-dimensional chiral algebra whose operator product expansion is
computed by soft-collinear limits. The comparison identifies this
algebra as the Mellin-shadow image of $\Phol_X$ applied to the
boundary soft-sector chiral algebra
$\cA^{\mathrm{soft}}_{\mathcal T}$ only after the celestial comparison
map
\[
 C_{\mathrm{cel},\mathcal T}:
 \Sh\Mel\bigl(\pi_r\Phol_X(\cA^{\mathrm{soft}}_{\mathcal T})\bigr)
 \longrightarrow \mathcal V^{\mathrm{cel}}_{\mathcal T}
\]
has been constructed and shown to preserve soft-collinear OPEs and
Ward identities.

\section{The $\SCchtop/E_3$ topologisation ladder}
\label{sec:part-iv-einfty-ladder}

The class-$M$ stratum of $\cM_{\ChirAlg}$ admits a canonical
ascending ladder in the coset slot of the fingerprint. The ladder
\begin{equation}
 \Vir_c \;\hookrightarrow\; \cW_N
 \;\hookrightarrow\; \cW_\infty[\mu]
 \;\hookrightarrow\; \cW_{1+\infty}[\psi]
 \;\xRightarrow{\;\SCchtop\;}\;
 \SCchtop(\cA_\infty,\Zder(\cA_\infty))
 \;\xRightarrow{\;T_{\mathrm{conf}},\,\widehat{\mathrm{wt}}\;}\;
 E_3^{\mathrm{top}}(\Zder(\cA_\infty)^\wedge)
 \label{eq:part-iv-einfty-ladder}
\end{equation}
has two distinct carriers. The $\SCchtop$ object is the generic
two-coloured chiral-topological structure on
$(\cA_\infty,\Zder(\cA_\infty))$. The $E_3^{\mathrm{top}}$ endpoint is
the additional topologisation obtained only after the conformal
vector, non-critical coupling, and weight-completed comparison
hypotheses are present.

\section{Moonshine as boundary-Hochschild $\Z/2$-orbifold}
\label{sec:part-iv-moonshine}

The Monster vertex operator algebra $V^\natural$
\cite{Frenkel-Lepowsky-Meurman} is the $\Z/2$-orbifold of the Leech
lattice vertex operator algebra $V_{\Lambda_{24}}$ by the involution
$\iota: v \mapsto -v$. The orbifold carrier is a
$\Z/2$-crossed VOA extension with anomaly class
$\omega_{\mathrm{DW}}\in H^3(\Z/2,U(1))\cong\Z/2$. The
Frenkel--Lepowsky--Meurman twisted-module construction supplies the
trivial associator class for this extension. A physical bulk reading
requires a separate open--closed map for the orbifold theory.

\section{Five Platonic theorems}
\label{sec:part-iv-five-theorems}

The five theorem surfaces below state the carrier-level content.  Each
theorem names its source coordinate, comparison map, and ambient
category; none identifies a physical object with a Vol I complex by
notation alone.

\subsection{Theorem 1: the associator carrier controls the named faces}
\label{sec:part-iv-thm-1}

\begin{theorem}[$\Face_K(\cA)$ carries a $\GRT_1(K)$ associator action;
 nine named comparison charts are typed; \ClaimStatusConditional]
\label{thm:part-iv-grt-torsor-of-bridges}
Let $\cA$ be a Koszul-compatible $\Eone$-chiral algebra in the
standard landscape, and let $K$ contain the coefficients of the
associators under consideration. The face scheme
$\Face_K(\cA)$ of Equation~\eqref{eq:part-iv-face-scheme} carries a
right action of $\GRT_1(K)$ induced by
Alekseev--Torossian--Kontsevich associator transport on the bar
differential. After passing to the fingerprint quotient
\[
 \Face_{K,\le 12}^{\mathrm{fp}}(\cA)
 :=
 \Face_K(\cA)\big/\bigl(F^{13}\mathfrak{grt}_1(K)+
 \ker(\mathrm{fp})\bigr),
\]
where $F^\bullet$ is the motivic weight filtration and
$\mathrm{fp}$ is the fingerprint coset map used in
Part~\ref{part:physics-bridges}, the following assertions hold.
\begin{enumerate}[label=\textup{(\roman*)}, leftmargin=*]
 \item $\Face_{K,\le 12}^{\mathrm{fp}}(\cA)$ is non-empty after a
 base associator in $\mathsf{Assoc}(K)$ is fixed;
 \item the nine named comparison charts
 \[
 F_1, F_2, F_3, F_4, F_5, F_6, F_7, F_8, F_9,
 \]
 have typed carriers: $F_1$ Vol~I bar hub; $F_2$
 Dimofte--Niu--Py line-operator $R$-twist; $F_3$ PVA
 $\lambda$-bracket r-matrix; $F_4$ Gaiotto--Zeng sphere
 Hamiltonians; $F_5$ Drinfeld Yangian on the finite-type affine
 lineage; $F_6$ Sklyanin--Semenov--Tian-Shansky classical
 r-matrix; $F_7$ Feigin--Frenkel--Reshetikhin Gaudin;
 $F_8$ Brown motivic MZV-coupled associator over the motivic-period
 field; $F_9$ Willwacher operadic chart through
 $H^0(\mathsf{GC}_2)\cong \mathfrak{grt}_1$;
 \item after a base associator is fixed, two charts $F_i,F_j$ are
 connected, when both charts are defined over the same field~$K$, by
 a class
 $\gamma_{ij}\in \GRT_1(K)/(F^{13}+\ker(\mathrm{fp}))$; the
 composite $\gamma_{ij} \circ \gamma_{jk} = \gamma_{ik}$ is the
 associator cocycle closure;
 \item no assertion of exhaustiveness of the full $\GRT_1(K)$-orbit
 follows from the nine named charts.  Exhaustiveness is a statement
 only about a specified finite fingerprint quotient and requires the
 quotient map $\mathrm{fp}$ as an input.
\end{enumerate}
\end{theorem}
\begin{proof}
Associator action: the bar differential of
$\bar{B}^{\mathrm{ch}}(\cA)$ is defined from the universal MC element
$\Theta_\cA$ via a Drinfeld associator~$\Phi$, and any two
associators over~$K$ differ by a unique element of~$\GRT_1(K)$
\cite{Drinfeld-Associator, Furusho-Associator}. The assignment
$\Phi \mapsto \bar{B}^{\mathrm{ch}}(\cA; \Phi)$ is equivariant; the
quotient by isomorphism gives the stated action on $\Face_K(\cA)$.
Non-emptiness follows from the chosen base associator. The named
charts are typed comparison maps out of the same collision residue:
the first seven are the bar, line-operator, PVA, Hamiltonian,
Yangian, Sklyanin--Semenov--Tian-Shansky, and Gaudin presentations;
the Brown chart lives over motivic periods rather than canonically
over~$\Q$; the Willwacher chart is defined over~$\Q$ at the
graph-complex Lie-algebra level.  Thus all nine are compatible only
after scalar extension to a common~$K$. Cocycle closure is the
composition law in the associator gauge groupoid.  Brown's motivic
freeness theorem \cite{Brown12} supplies the odd-weight motivic generators; it does
not by itself prove that a finite list of nine charts exhausts the
full $\GRT_1(K)$-orbit.  The finite-atlas assertion is therefore
exactly the assertion of the chosen fingerprint quotient.
$\quad\square$
\end{proof}

\begin{remark}[Coefficient fields for the two new faces]
\label{rem:part-iv-thm-1-hziv}
The Brown face is motivic-period valued: the KZ associator has
multiple-zeta coefficients, and Brown's theorem controls their
motivic coaction.  It is not a canonical $\Q$-point of
$\GRT_1$.  The Willwacher face is rational at the level of the
graph-complex theorem
$H^0(\mathsf{GC}_2)\cong\mathfrak{grt}_1$.  Comparing $F_8$ with
$F_9$ therefore requires scalar extension to a common coefficient
field and a chosen realisation functor from motivic periods.
\end{remark}

\subsection{Theorem 2: coordinate projections and comparison maps}
\label{sec:part-iv-thm-2}

\begin{theorem}[The boundary-holographic functor supplies the Vol I
 source coordinates; physical bridges require comparison maps;
 \ClaimStatusConditional]
\label{thm:part-iv-universal-holography-unifies}
For every Koszul-compatible $\Eone$-chiral algebra~$\cA$ in the
standard landscape, the functor
$\Phol_X: \mathsf{ChirAlg}^{\mathrm{Kosz}}_X \to \Holbd_X$
of Definition~\ref{def:part-iv-Phol} admits five canonical
projection functors
\begin{equation}
 \pi_{\mathrm{bar}},\;\; \pi_{\MC},\;\; \pi_{r},\;\;
 \pi_{\mathrm{Hoch}},\;\; \pi_{\nabla},
\end{equation}
onto respectively the bar coalgebra, the Maurer--Cartan moduli, the
classical r-matrix hub, the derived chiral centre, and the
Gauss--Manin connection. The physics-facing identifications factor
through the following named comparison maps when those maps exist:
\begin{enumerate}[label=\textup{(\roman*)}, leftmargin=*]
 \item the Feynman comparison is a map from
 $\pi_{\MC}\Phol_X(\cA)$ to the worldline classical action;
 \item the BV--BRST comparison is the open--closed map
 $\OC_\cA$ from the physical BRST closed sector to
 $\pi_{\mathrm{Hoch}}\Phol_X(\cA)=\Zder(\cA)$;
 \item the derived-Langlands comparison is the critical-level
 specialization of $\pi_\nabla\Phol_X(\cA)$ at $k=-h^\vee$ to the
 oper connection;
 \item the arithmetic-shadows comparison pairs
 $\pi_{\mathrm{Hoch}}\Phol_X(\cA)$ with the Selberg/Eisenstein
 regularization after the relevant analytic completion has been
 supplied;
 \item the face comparison is the map from
 $\pi_r\Phol_X(\cA)$ to the named charts of
 $\Face_K(\cA)$ in Theorem~\ref{thm:part-iv-grt-torsor-of-bridges}.
\end{enumerate}
Thus the Vol I theorem supplies the five source coordinates.  A
physics bridge is an equality only after its comparison map is a
quasi-isomorphism, an isomorphism of factorisation algebras, or a
specified equality in the target category.
\end{theorem}
\begin{proof}
Existence of $\Phol_X$: Definition~\ref{def:part-iv-Phol} assembles
the five coordinates; functoriality follows from functoriality of
the bar coalgebra
\textup{(}Theorem~\ref{thm:operadic-homotopy-convolution}\textup{)},
of the universal MC element
\textup{(}Theorem~\ref{thm:mc2-bar-intrinsic}\textup{)}, and of the
chiral Hochschild complex
\textup{(}Chapter~\ref{ch:en-koszul-duality}\textup{)}. The
projections $\pi_{\mathrm{bar}}$ etc.\ are forgetful functors
onto the indicated factor of the boundary-holographic datum; each is
well-defined on morphisms.

The five factorizations of the Part~\ref{part:physics-bridges}
chapters are comparison statements, not definitions of $\Phol_X$.
The Feynman statement uses the bar-worldline dictionary
$\Theta_\cA=S_{\mathrm{cl}}$ on its proved surface
\textup{(}Theorem~\ref{thm:brst-bar-genus0}\textup{)}.  The BV--BRST
statement uses the coderived comparison
\textup{(}Theorem~\ref{thm:bv-bar-coderived-vol1}\textup{)} and is
exactly an $\OC_\cA$ statement.  The critical-level statement uses
Theorem~\ref{thm:oper-bar-dl}.  The arithmetic statement requires the
Selberg regularization in Chapter~\ref{chap:arithmetic-shadows}.  The
face statement is Theorem~\ref{thm:part-iv-grt-torsor-of-bridges}.
Forgetting the comparison maps leaves only the five source
coordinates of $\Phol_X(\cA)$, so no physical equality is asserted
without its carrier. $\quad\square$
\end{proof}

\begin{remark}[The derived-centre coordinate]
\label{rem:part-iv-thm-2-hziv}
The coordinate $\pi_{\mathrm{Hoch}}\Phol_X(\cA)$ is
$\Zder(\cA)=C^\bullet_{\mathrm{ch}}(\cA,\cA)$. It is a brace dg
algebra in the factorisation bimodule category. The phrase
``physical bulk'' denotes this complex only on a field-theoretic
surface where $\OC_\cA$ has been constructed.
\end{remark}

\subsection{Theorem 3: celestial factorization comparison}
\label{sec:part-iv-thm-3}

\begin{theorem}[Celestial factorization is the Mellin-shadow image of
 the boundary source under the soft comparison map;
 \ClaimStatusConditional.  The celestial algebra construction is
 proved elsewhere by Pasterski--Shao--Strominger and Strominger; the
 class-$M$ chain-level comparison at $g \ge 1$ uses Vol~II
 \texttt{V2-thm:uch-main} and the hypotheses of
 \texttt{V2-conj:uch-gravity-chain-level}.]
\label{thm:part-iv-celestial-is-holographic-image}
Let $\mathcal{T}$ be a four-dimensional gauge theory in its soft
sector with asymptotic null infinity~$\mathscr{I}^\pm$, and let
$\cA^{\mathrm{soft}}_{\mathcal{T}}$ denote the chiral soft algebra
generated by its leading soft-theorem operators. Assume:
\begin{enumerate}[label=\textup{(\alph*)}, leftmargin=*]
 \item the soft-sector open--closed map
 $\OC_{\mathrm{soft},\mathcal T}:
 \mathcal O^{\mathrm{phys}}_{\mathrm{soft}}(\mathcal T)
 \to \Zder(\cA^{\mathrm{soft}}_{\mathcal T})$ is a
 quasi-isomorphism on the operators used in the celestial OPE;
 \item the residue comparison sends the soft-collinear leading pole
 to $r_{\cA^{\mathrm{soft}}_{\mathcal T}}(z)
 =\Res^{\mathrm{coll}}_{0,2}(\Theta_{\cA^{\mathrm{soft}}_{\mathcal T}})$;
 \item the Mellin and shadow transforms are taken in the
 Pasterski--Shao--Strominger normalization.
\end{enumerate}
Then:
\begin{enumerate}[label=\textup{(\roman*)}, leftmargin=*]
 \item the Mellin transform
 $\Mel: \cA^{\mathrm{soft}}_{\mathcal{T}} \to
 \cA^{\mathrm{soft}}_{\mathcal{T}}(\Delta)$
 converts energy eigenstates into conformal-dimension-$\Delta$
 primaries;
 \item the shadow transform $\Sh: \Delta \mapsto 2 - \Delta$
 produces the shadow-primary basis adapted to two-dimensional
 chiral factorization on the celestial sphere~$\CP^1$;
 \item the composite $\Sh \circ \Mel \circ \Phol_X:
 \mathsf{ChirAlg}^{\mathrm{Kosz}}_{X} \to
 \FactAlg_{\CP^1}$ lands in the category of holomorphic
 factorization algebras on the celestial sphere;
 \item the celestial factorization algebra
 $\mathcal{V}^{\mathrm{cel}}_{\mathcal{T}}$ of $\mathcal{T}$ is
 identified with the image
 $\Sh(\Mel(\Phol_X(\cA^{\mathrm{soft}}_{\mathcal{T}})))$ by the
 comparison map $C_{\mathrm{cel},\mathcal T}$;
 \item the operator product expansion of
 $\mathcal{V}^{\mathrm{cel}}_{\mathcal{T}}$ is computed by
 soft-collinear limits, and these limits are the Mellin-shadow
 image of the classical r-matrix $r(z)$-coordinate of
 $\Phol_X$.
\end{enumerate}
\end{theorem}
\begin{proof}
Mellin transform: the standard Mellin transform on the
positive-energy half of the asymptotic momentum
space~$\mathbb{R}_{> 0}$ is well-defined on soft eigenmodes
\cite{Pasterski-Shao-Strominger-Mellin}; it sends energy
eigenstates to conformal primaries of boost weight~$\Delta$.
Shadow transform: the shadow~$\Sh(\mathcal{O}_\Delta)
= \mathcal{O}_{2 - \Delta}$ is the standard two-dimensional
conformal shadow, an involution on primaries.
Composite: the functor $\Phol_X$ produces a chiral coalgebra
equipped with collision-residue data; Mellin transform converts
this chiral coalgebra to a celestial-sphere $\Eone$-coalgebra;
shadow transform relates the $\Delta$-lift to its $2-\Delta$
shadow, recovering holomorphic factorization on the celestial
sphere.

Celestial comparison: the celestial factorization algebra
$\mathcal{V}^{\mathrm{cel}}_{\mathcal{T}}$ is constructed in the
literature from the soft-theorem operators and the
boost-eigenstate basis \cite{Strominger-Celestial,
Pasterski-Shao-Strominger-Mellin}; this construction is
manifestly the composite of (Mellin) followed by (shadow transform)
applied to the chiral soft algebra.  The chiral soft algebra is the
image of~$\Phol_X$ restricted to the four-dimensional soft sector
only under the assumed open--closed soft comparison.  Residue
compatibility identifies the soft-theorem operators with the collision
residues of the classical r-matrix in the soft-limit normalization.
Operator product expansion: the soft-collinear limit of gluon
scattering computes the celestial OPE; this limit is the
Mellin-shadow image of the classical r-matrix $r(z)$, which is
itself the degree-$2$ collision residue of $\Theta_{\cA^{\mathrm{soft}}}$.
$\quad\square$
\end{proof}

\begin{remark}[Celestial OPE test]
\label{rem:part-iv-thm-3-hziv}
The comparison map $C_{\mathrm{cel},\mathcal T}$ is tested on the
leading soft-gluon OPE computed from soft-collinear limits
\cite{Fan-Fotopoulos-Taylor-Celestial, Strominger-Celestial}.  In the
gravitational soft sector the same test is the
$\cW_{1+\infty}$ Ward algebra of Strominger
\cite{Strominger-w1plusinfty}.  These computations check the
$r_\cA(z)$ coordinate; they do not identify a physical bulk with
$\Zder(\cA)$ without $\OC_{\mathrm{soft},\mathcal T}$.
\end{remark}

\subsection{Theorem 4: The $\SCchtop/E_3$ topologisation ladder}
\label{sec:part-iv-thm-4}

\begin{theorem}[Class-$M$ ladder through $\SCchtop$ and, under
 topologisation hypotheses, to $E_3^{\mathrm{top}}$;
 \ClaimStatusConditional\ on the weight-completed ambient;
 \ClaimStatusProvedElsewhere\ for the Vol~II endpoint established by
 \texttt{V2-thm:e-infinity-topologization-ladder} and
 \texttt{V2-thm:iterated-sugawara-construction}. On the
 bounded-direct-sum ambient $\mathrm{Ch}(\mathrm{Vect})$ the
 chain-level statement is genuinely false (see
 Theorem~\ref{thm:topologization-tower}); the statement below is
 the weight-completed reformulation]
\label{thm:part-iv-e-infinity-ladder}
Let $(\cA_N)_{N \ge 2}$ be the principal $\cW_N$-tower in the
class-$M$ stratum of $\cM_{\ChirAlg}$, with
$\cA_2 = \Vir_c$ the Virasoro algebra, $\cA_N = \cW_N$ the
principal $\cW$-algebra of $\fsl_N$-type, and
$\cA_\infty = \cW_{1+\infty}[\mu]$ the two-parameter
$\cW$-infinity algebra. The ladder of inclusions
\begin{equation}
 \cA_2 \;\hookrightarrow\; \cA_3 \;\hookrightarrow\; \cdots
 \;\hookrightarrow\; \cA_\infty
\end{equation}
admits the following two-stage carrier:
\begin{enumerate}[label=\textup{(\roman*)}, leftmargin=*]
 \item each inclusion $\cA_N \hookrightarrow \cA_{N+1}$ preserves
 the Koszul conductor
 $\kappa_{\mathrm{ch}}(\cA_N) = c\cdot(H_N - 1)$,
 with $H_N = \sum_{j=1}^{N} 1/j$
 \textup{(}by the universal formula of
 Theorem~\ref{thm:part-iii-central-charge-is-emergent} at
 the class-$M$ stratum\textup{)};
 \item the direct limit $\cA_\infty = \mathrm{colim}_N \cA_N =
 \cW_{1+\infty}[\mu]$ is a two-parameter $\cW$-algebra whose
 Koszul conductor diverges as $H_\infty$;
 \item the pair
 $(\cA_\infty,\Zder(\cA_\infty))$ carries the generic
 $\SCchtop$ boundary/closed-sector structure;
 \item at non-critical Sugawara coupling the class-$M$
 topologization theorem
 \textup{(}Theorem~\ref{thm:topologization-tower}\textup{)}
 gives, on a strict pro-object / $J$-adic / filtered
 weight-completed presentation satisfying the continuous comparison
 hypotheses, an $E_3^{\mathrm{top}}$
 structure on the weight-completed chiral Hochschild complex
 $\mathrm{ChirHoch}^\bullet(\cA_\infty)^\wedge$; chain-level closure
 on the bounded-direct-sum ambient is genuinely false for class $M$
 and is not claimed;
 \item the $\GRT_1(K)$ associator action transports the
 $\SCchtop$ carrier; it descends to the $E_3^{\mathrm{top}}$
 multiplication exactly when the topologisation map
 $\SCchtop(\cA_\infty,\Zder(\cA_\infty))\to
 E_3^{\mathrm{top}}(\Zder(\cA_\infty)^\wedge)$ is
 $\GRT_1(K)$-equivariant.
\end{enumerate}
\end{theorem}
\begin{proof}
Tower structure: the principal $\cW$-algebras are nested by
Feigin--Frenkel screening: $\cW_N$ is recovered inside $\cW_{N+1}$
as the kernel of an additional screening charge
\cite{Fateev-Lukyanov}, so the inclusions are maps of vertex
algebras. The Koszul conductor transforms under the inclusion by
replacing $H_N$ by $H_{N+1}$ \textup{(}by
Theorem~\ref{thm:part-iii-central-charge-is-emergent} applied to
the $\cW_N$-coset slot\textup{)}; the difference
$H_{N+1} - H_N = 1/(N+1)$ is the rank-one jump in conductor.
Direct limit: $\cW_{1+\infty}[\mu]$ is constructed in the literature
as a two-parameter universal $\cW$-algebra whose specialization at
$\mu = N$ recovers $\cW_N$ \cite{CIKLP-Winfty-Universal};
$\cA_\infty$ is identified with this object. Conductor divergence:
$\kappa_{\mathrm{ch}}(\cA_\infty) = c \cdot H_\infty$ diverges as a
harmonic series; the class-$M$ topologization is restricted to the
weight-completed Hochschild complex where this divergence is
regulated by the weight filtration.

The $\SCchtop$ carrier is the two-coloured chiral-topological
structure on the boundary algebra and its derived centre.  The
topologization at non-critical Sugawara is stronger: the $E_3^{\mathrm{top}}$
structure on $\mathrm{ChirHoch}^\bullet(\cA_\infty)^\wedge$
follows from the class-$M$ weight-completed topologization theorem
\textup{(}Theorem~\ref{thm:topologization-tower}, class-$M$
direction, weight-completed\textup{)}; the
$E_3$ multiplication is defined on the completed comparison surface
and passes to the completed direct limit by the
Milnor--Mittag-Leffler filtration argument.

$\GRT_1(K)$-equivariance: the $\cW_{1+\infty}[\mu]$-algebra carries an
associator action on the coefficient field~$K$
\cite{Gaiotto-Rapcak-Y}.  This action is automatically defined on the
$\SCchtop$ carrier.  Descent to $E_3^{\mathrm{top}}$ is the
equivariance of the topologisation map, not a formal consequence of
Dunn additivity in the coloured operad. $\quad\square$
\end{proof}

\begin{remark}[Class-$M$ ambient]
\label{rem:part-iv-thm-4-hziv}
For class $M$, the bounded direct-sum chain complex does not carry
the asserted chain-level topologisation.  The true endpoint lies in
the pro-object / $J$-adic / filtered weight-completed ambient.  The
Linshaw--Gaiotto--Rapcak construction of
$\cW_{1+\infty}[\mu]$ supplies the infinite-rank algebraic carrier;
the Costello--Gwilliam holomorphic-topological factorization algebra
supplies an independent $E_3^{\mathrm{top}}$ model only on the
comparison surface where the topologisation map exists.
\end{remark}

\subsection{Theorem 5: Moonshine as boundary-Hochschild orbifold}
\label{sec:part-iv-thm-5}

\begin{theorem}[Monster $V^\natural$ as boundary-Hochschild
 $\Z/2$-orbifold carrier; Dijkgraaf--Witten class vanishes for the
 FLM extension; external FLM--Borcherds--Gannon inputs]
\label{thm:part-iv-moonshine-via-holography}
\ClaimStatusConditional
The FLM orbifold construction of $V^\natural$, the Monster
denominator identity, and the genus-zero theorem for McKay--Thompson
series are external inputs.
The Monster vertex operator algebra
$V^\natural$ \cite{Frenkel-Lepowsky-Meurman} is the
boundary-Hochschild $\Z/2$-orbifold carrier of the Leech lattice vertex operator
algebra $V_{\Lambda_{24}}$ under the involution
$\iota: v \mapsto -v$:
\begin{enumerate}[label=\textup{(\roman*)}, leftmargin=*]
 \item the FLM extension is
 $V^\natural=V_{\Lambda_{24}}^+\oplus V_{\Lambda_{24}}^{T,+}$;
 \item this holomorphic extension is obstructed by an anomaly
 class
 $\omega_{\mathrm{DW}} \in H^3(\Z/2, U(1)) \cong \Z/2$;
 \item $\omega_{\mathrm{DW}} = 0$ on the $V_{\Lambda_{24}}$
 orbifold; the vanishing is the associator trivialization in the
 Frenkel--Lepowsky--Meurman $\iota$-twisted-module construction,
 while
 $\Theta_{\Lambda_{24}}(\tau)/\eta(\tau)^{24}=J(\tau)+24$
 verifies the Leech graded-character shadow;
 \item the resulting $\Etwo$-topological structure on
 $\Zder(V^\natural) = \mathrm{ChirHoch}^\bullet(V^\natural)$ is the
 orbifold-Hochschild carrier attached to the moonshine module, not a
 physical bulk before an open--closed map is supplied;
 \item the character-level constants are
 $\mathrm{ch}(V_{\Lambda_{24}})=J+24$ and
 $\mathrm{ch}(V^\natural)=J$, equivalently
 $(V^\natural)_0=\mathbb C\mathbf 1$ and $(V^\natural)_1=0$.
\end{enumerate}
\end{theorem}
\begin{proof}
Construction of $V^\natural$ as orbifold: Frenkel--Lepowsky--Meurman
construct $V^\natural$ as the $\Z/2$-orbifold of
$V_{\Lambda_{24}}$ by the involution $\iota: v \mapsto -v$
\cite{Frenkel-Lepowsky-Meurman}; the orbifold has vacuum
character $\chi_0(\tau) = J(\tau)$, where $J=j-744$.

$\Etwo$-topological structure on the orbifold-Hochschild carrier:
the $\Z/2$-orbifold
of a chiral algebra~$\cA$ extends to an orbifold of its chiral
Hochschild complex $\mathrm{ChirHoch}^\bullet(\cA)$ if and only if
the anomaly class in $H^3(\Z/2, U(1))$ vanishes
\cite{Dijkgraaf-Witten, Kirillov-Orbifold}; when the class vanishes,
the orbifold carries a compatible $\Etwo$-topological multiplication
\textup{(}by the boundary-Hochschild topologization argument of
Chapter~\ref{chap:topologization-chain-level-platonic}
applied to the orbifold-Hochschild complex\textup{)}.

Vanishing of $\omega_{\mathrm{DW}}$ for $V_{\Lambda_{24}}$: the
Dijkgraaf--Witten class of the $\Z/2$-orbifold on
$V_{\Lambda_{24}}$ is the associator obstruction for the
$\Z/2$-crossed extension
\[
 V_{\Lambda_{24}} \oplus V_{\Lambda_{24}}^{T,+}
\]
by the even part of the canonical $\iota$-twisted module. In the
Frenkel--Lepowsky--Meurman construction the central extension
$1\to\{\pm1\}\to\widehat{\Lambda}_{24}\to\Lambda_{24}\to1$ has
commutator $(-1)^{\langle\alpha,\beta\rangle}$, and the chosen
normalized $2$-cocycle has trivial coboundary on the even unimodular
Leech lattice. The twisted Jacobi identity for
$V_{\Lambda_{24}}^{T}$ is therefore an honest associativity identity,
not a projective identity. Under the standard identification between
$H^3(\Z/2,U(1))$ and Dijkgraaf--Witten associator phases, this is
exactly $\omega_{\mathrm{DW}}=0$. The theta-series identity
$\Theta_{\Lambda_{24}}(\tau)/\eta(\tau)^{24} = J(\tau)+24$
\cite{Conway-Sloane, Borcherds-Monster} and the Monster denominator
identity verify the graded character and denominator of the resulting
orbifold; they are character-level checks of the already constructed
cocycle trivialization.

Moonshine module identification: with $\omega_{\mathrm{DW}} = 0$, the
orbifold admits a consistent extension; the
Conway--Norton--Borcherds moonshine module~$V^\natural$
\cite{Borcherds-Monster} is this extension. The Monster symmetry is
the automorphism group of the resulting VOA; it is not determined by
the vanishing of the Dijkgraaf--Witten class alone.

Character constants: the Leech lattice character is
$\Theta_{\Lambda_{24}}/\eta^{24}=J+24$; the FLM orbifold cancels the
weight-one Leech contribution and gives $\mathrm{ch}(V^\natural)=J$.
Thus $(V^\natural)_0=\mathbb C\mathbf 1$ and
$(V^\natural)_1=0$.  These constants check the VOA orbifold; they are
not classes in $H^3(\Z/2,U(1))$. $\quad\square$
\end{proof}

\begin{remark}[Moonshine checks outside Hochschild theory]
\label{rem:part-iv-thm-5-hziv}
Borcherds' proof of monstrous moonshine uses the Monster denominator
identity, and Gannon's genus-zero theorem uses modular functions.
Those results identify the moonshine module and its McKay--Thompson
series independently of the chiral Hochschild carrier.  They do not
replace the Dijkgraaf--Witten obstruction calculation for the
$\Z/2$-crossed extension.
\end{remark}

\section{Projection algebra}
\label{sec:part-iv-reading-order}

Let
\[
 \Pi=\{\pi_{\mathrm{bar}},\pi_{\MC},\pi_r,\pi_{\mathrm{Hoch}},
 \pi_\nabla\}
\]
be the five projections of $\Phol_X(\cA)$. A physics bridge with
target object $T_\beta$ is not the projection $\pi_i\Phol_X(\cA)$
alone. It is a pair
\[
 \bigl(\pi_i,\; C_\beta:
 F_\beta(\pi_i\Phol_X(\cA))\longrightarrow T_\beta\bigr),
\]
where $F_\beta$ is the required normalisation, completion, Mellin
transform, shadow transform, critical-level specialization, or
orbifold construction. The bridge theorem is the assertion that
$C_\beta$ is an isomorphism, quasi-isomorphism, or equality in its
declared target category.

The face scheme is the same rule in the associator direction:
$F_i$ is a comparison chart for $r_\cA(z)$ inside
$\Face_K(\cA)$, not a new object replacing $\Phol_X(\cA)$.

\section{Comparison carriers}
\label{sec:part-iv-dependency-map}

The five carrier statements
(\ref{thm:part-iv-grt-torsor-of-bridges},
\ref{thm:part-iv-universal-holography-unifies},
\ref{thm:part-iv-celestial-is-holographic-image},
\ref{thm:part-iv-e-infinity-ladder},
\ref{thm:part-iv-moonshine-via-holography}) have the following
source coordinates and comparison carriers.

\begin{center}
\begin{tabular}{p{3.8cm}|p{6.2cm}|p{4.5cm}}
\hline
\textbf{Statement} & \textbf{Vol I source coordinate}
& \textbf{comparison carrier} \\
\hline
\ref{thm:part-iv-grt-torsor-of-bridges}
 ($\Face_K(\cA)$)
 & Vol~II Ch.~\ref{V2-ch:grt-seven-faces}
 & associator torsor over $K$; fingerprint quotient map \\
\hline
\ref{thm:part-iv-universal-holography-unifies}
 ($\Phol_X$ projections)
 & Ch.~\ref{ch:v1-feynman},
 Ch.~\ref{ch:v1-bv-brst},
 Ch.~\ref{ch:derived-langlands},
 Ch.~\ref{chap:arithmetic-shadows}
 & comparison maps $C_\beta$ and $\OC_\cA$ where physical bulk is
 named \\
\hline
\ref{thm:part-iv-celestial-is-holographic-image}
 (celestial comparison)
 & Ch.~\ref{ch:v1-feynman}
 & $C_{\mathrm{cel},\mathcal T}$ after Mellin, shadow, and
 soft-sector open--closed comparison \\
\hline
\ref{thm:part-iv-e-infinity-ladder}
 ($\SCchtop/E_3$ topologisation)
 & Ch.~\ref{chap:w-algebras},
 Ch.~\ref{chap:topologization-chain-level-platonic}
 & $\SCchtop(\cA,\Zder(\cA))\to
 E_3^{\mathrm{top}}(\Zder(\cA)^\wedge)$ under completion and
 conformal-vector hypotheses \\
\hline
\ref{thm:part-iv-moonshine-via-holography}
 (Monster orbifold carrier)
 & Ch.~\ref{chap:chiral-moonshine-unified},
 Ch.~\ref{ch:moonshine}
 & FLM $\Z/2$-crossed extension; Dijkgraaf--Witten class in
 $H^3(\Z/2,U(1))$ \\
\hline
\end{tabular}
\end{center}

\section{Carrier theorem}
\label{sec:part-iv-bridges-restated}

\begin{proposition}[Part~\ref{part:physics-bridges} restated as
 carrier data with typed comparison maps;
 \ClaimStatusConditional]
\label{prop:part-iv-bridges-restated}
Under this structural reading, the content of
Part~\ref{part:physics-bridges} is the datum
\begin{equation}
 \bigl(\Phol_X,\; \Face_K(\cA),\;
 \{C_\beta\}_{\beta\in\{\mathrm{Feyn},\mathrm{BRST},\mathrm{DL},
 \mathrm{arith},\mathrm{cel},\mathrm{orb},\mathrm{top}\}}\bigr),
 \label{eq:part-iv-bridges-triple}
\end{equation}
where $\Phol_X$ is Definition~\ref{def:part-iv-Phol},
$\Face_K(\cA)$ is Equation~\eqref{eq:part-iv-face-scheme}, and each
$C_\beta$ is a comparison map in the target category of the
corresponding carrier.  The celestial algebra, the Monster orbifold,
the K3/BKM arithmetic surface, and the $\SCchtop/E_3$ topologisation
are not Vol I theorems until their comparison maps are part of the
statement.
\end{proposition}
\begin{proof}
The projections of $\Phol_X$ are functorial by
Theorem~\ref{thm:part-iv-universal-holography-unifies}.  The
associator carrier is Theorem~\ref{thm:part-iv-grt-torsor-of-bridges}.
The celestial, topologisation, and Monster assertions are precisely
Theorems~\ref{thm:part-iv-celestial-is-holographic-image},
\ref{thm:part-iv-e-infinity-ladder}, and
\ref{thm:part-iv-moonshine-via-holography}, each with its named
comparison carrier.  The datum \eqref{eq:part-iv-bridges-triple}
therefore records exactly the maps needed to turn a Vol I coordinate
into a physics-facing object. $\quad\square$
\end{proof}

\medskip

Every bridge in Part~\ref{part:physics-bridges} is a coordinate of
$\Phol_X$ together with a comparison map.  The comparison map is the
mathematical place where physical content enters.

\section{The class-$\mathcal M$ vs class-$\mathcal C$ dichotomy and its $\mathbb Z/8$ shared core}
\label{sec:pivpi-supplement}

The K3/BKM material in this section is a comparison carrier for
Vol~III, not a Vol I theorem identifying the bar complex with the
Hall--Borcherds object.  The carrier is the map from the curve-level
boundary datum to the Hall--Borcherds surface
\[
 \mathbf H_{\Delta_5}\leadsto
 \mathcal D_\hbar\!\left(
 \mathcal Y^{\mathrm{Hall}}_\hbar(\mathrm{CoHA}_{K3\times E})\right),
\]
together with the period, Humbert-monodromy, and Dijkgraaf--Witten
cocycle comparisons stated below.  Equalities such as
$K^{\kappa_{\mathrm{ch}}}=8$ are normalisations on this carrier until
the comparison map is part of the theorem.

\begin{remark}[The $\mathcal M$ vs $\mathcal C$ dichotomy]
\label{rem:pivpi-1}
The K3/BKM comparison refines the shadow-tower quadrichotomy into two
major arithmetic classes: $\mathcal M$ (Mathieu; K3, $M_{24}$, $\Delta_5$)
and $\mathcal C$ (Conway; Leech, $\mathrm{Co}_0$, Borcherds Monster).
The dichotomy meets at the $\mathbb Z/8$-class shared between Mukai
doubling on K3 ($K^{\kappa_{\mathrm{ch}}} = 8$) and the Leech
$(-1)$-automorphism shift on $\mathrm{Co}_0$ (order-$8$ in the centre)
\cite{ConwaySloane1988,Borcherds1998}. The Vol.~III carrier is
Chapter~\ref{V3-ch:k3-chiral-bialgebra-platonic}.
\end{remark}

\begin{remark}[Humbert $H_4$ monodromy and the
$\overline{\mathcal A_2} \setminus (H_1 \cup H_4)$ fundamental group]
\label{rem:pivpi-humbert-H4}
The Koszul locus of the $\mathcal M$-family lives on
$\overline{\mathcal A_2} \setminus (H_1 \cup H_4)$, the compactified
Siegel threefold with both the first Humbert divisor $H_1$
(discriminant-$1$ real multiplication, product locus
$\mathcal A_1 \times \mathcal A_1$) and the fourth Humbert divisor
$H_4$ (discriminant-$4$ real multiplication, squared-CM locus) removed.
The local Gauss--Manin monodromy of the holonomic $D$-module
$\mathcal L^{\Delta_5}$ has order $8$ at $H_1$ and order $16$ at
$H_4$ by Bruinier's Heegner-divisor reciprocity, in the form recorded
as Vol~III Theorem~\ref{V3-thm:humbert-order-K-kappa}. The
discriminant-$4$ CM-degeneration of the genus-$2$ period carries a
Kummer parity quotient of order $2$
\cite{HulekSankaran2002,vdGeer1988}. Locally around the two divisors,
$\pi_1^{\mathrm{loc}}(\overline{\mathcal A_2} \setminus
(H_1 \cup H_4))$ maps onto cyclic monodromy quotients
$\Z/8$ and $\Z/16$; its Kummer quotient replaces the second factor by
$\Z/2$. The two local loops commute only up to the Mayer--Vietoris
boundary map on $H_1 \cap H_4 \subset \overline{\mathcal A_2}$; at
smooth transversal strata the commuting monodromy quotient is
$\Z/8 \times \Z/16$, while the parity quotient is
$\Z/8 \times \Z/2$. The $\Z/8$ Mukai-doubling identity
$K^{\kappa_{\mathrm{ch}}} = 8$ (Remark~\ref{rem:pivpi-1}) accounts
for the $H_1$ factor; the $M_{24}$ order-$6$ sector below detects the
Kummer parity quotient of the $H_4$ monodromy alone. Primary:
Hulek--Sankaran 2002 ``The geometry of Siegel modular varieties''
(Adv.\ Stud.\ Pure Math.\ 35); Bruinier 2002 Prop.\ 5.1 (Invent.\ Math.);
van der Geer 1988 ``Hilbert modular surfaces'' (Springer). The
corresponding Vol.~III carrier is
Chapter~\ref{V3-ch:k3-chiral-bialgebra-platonic}.
\end{remark}

\begin{remark}[Kazhdan--Lusztig equivalence scope refined across $H_1 \cup H_4$]
\label{rem:pivpi-KL-scope}
The chain-level Kazhdan--Lusztig equivalence
$\mathrm{Mod}(\mathbf H_{\Delta_5})^{\mathrm{loc,KL}} \simeq
\mathrm{Mod}(\mathfrak u_{q_{\mathrm{K3}}}(\widehat{\mathfrak m}_{\mathrm{K3}}))$
at $q = e^{2\pi i / 8}$ is proved \emph{on the Koszul locus}
$\overline{\mathcal A_2} \setminus (H_1 \cup H_4)$. Its extension
across each stratum requires a different Positselski regime:
\begin{itemize}[leftmargin=1.5em,itemsep=0.25em]
 \item Across $H_1$: use the weight-completed coderived/contraderived
 model (Positselski 2011, LNM 1780, Ch.\ 6), pulled back along the
 $\Z/8$ local monodromy. The Positselski comparison gives an
 isomorphism at the weight-completed level and an obstruction class
 in the $\Z/8$-equivariant cohomology of $H_1$ that measures the
 failure of strict equality.
 \item Across $H_4$: use the $\Z/16$-monodromic KL regime together
 with its Kummer parity quotient. The order-$2$ quotient lifts to a
 Kummer involution on the chiral Hall algebra, and the KL equivalence
 extends through the resulting graded bimodules; the Kummer
 specialisation (Bruinier LNM 1780 \S 4) identifies the $(+)$-eigenspace
 with the K3 side and the $(-)$-eigenspace with the twisted-K3 side.
 \item On the Koszul locus: the chain-level statement with
 $q_{\mathrm{K3}} = e^{2\pi i / 8}$ holds strictly, no
 weight-completion required.
\end{itemize}
Primary: Positselski 2011 ``Two kinds of derived categories, Koszul
duality, and comodule-contramodule correspondence'' (Mem.\ AMS LNM
1780, via arXiv:0905.2621); Bruinier LNM 1780; Kazhdan--Lusztig 1993
I--IV (J.\ AMS). The compatible carriers are Vol.~III
Chapter~\ref{V3-ch:k3-chiral-bialgebra-platonic} and Vol.~II
Chapter~\ref{chap:periodic-cdg-admissible}.
\end{remark}

\begin{remark}[$6A$ vs $6B$ class distinguishes the order-$6$ sector of the $\mathcal M$ apex]
\label{rem:pivpi-M24-order6}
Inside the class-$\mathcal M$ apex $\mathbf H_{\Delta_5}$, the
$M_{24}$-action splits the order-$6$ sector into two pieces by
conjugacy class. The $6A$ class with cycle shape $1^2 2^2 3^2 6^2$
and umbral shift $m_{6A} = 2$ contributes to the Humbert-$H_4$
parity quotient via its pure $\Z/3$-phase (CRT residue $(0, 2)$
in $\Z/6 = \Z/2 \times \Z/3$, pentagon-$\hbar^3$ cocycle class
$2 \pmod 6$). The $6B$ class with cycle shape $6^4$ and umbral shift
$m_{6B} = 6$ contributes the $\Z/2$ sign-flip of the Humbert $H_4$
monodromy via its pure $\Z/2$-phase (CRT residue $(1, 0)$, cocycle
class $3 \pmod 6$) \cite{ChengDuncanHarvey2014,ConwayCurtisNortonParkerWilson1985}.
Their sum $2 + 3 \equiv -1 \pmod 6$ is the order-$6$ restriction of
the universal identity $\hbar^2 \cdot K^{\kappa_{\mathrm{ch}}} = -1$.
This is the cyclic restriction of Vol~III's pentagon--umbral agreement
theorem and Vol.~II Remark~\ref{V2-rem:was-6A-6B}.
\end{remark}

\begin{remark}[$v_{\Delta_5}$ character table on $\mathrm{Sp}_4(\Z)$ generators]
\label{rem:pivpi-vDelta5-chartable}
The sign character
$v_{\Delta_5} \colon \mathrm{Sp}_4(\Z) \to \{\pm 1\}$ factors through
$\mathrm{Sp}_4(\Z/2) \cong S_6$ (Igusa $1964$, through the action on
theta characteristics); the kernel
$\ker(v_{\Delta_5}) \subset \mathrm{Sp}_4(\Z)$ is the index-$2$
subgroup of matrices mapping to even permutations in $S_6$. On the
Klingen--Igusa generators
\[
 J = \begin{pmatrix} 0_2 & -I_2 \\ I_2 & 0_2 \end{pmatrix},\qquad
 T_b = \begin{pmatrix} I_2 & b \\ 0_2 & I_2 \end{pmatrix}\ (b^t = b),\qquad
 M_U = \begin{pmatrix} U & 0 \\ 0 & (U^t)^{-1} \end{pmatrix}\ (U \in \mathrm{GL}_2(\Z)),
\]
the character values are
\[
 v_{\Delta_5}(J) = +1,\qquad
 v_{\Delta_5}(T_b) = (-1)^{b_1+b_2+b_3},\qquad
 v_{\Delta_5}(M_U)
 = (-1)^{(1+u_1+u_4)(1+u_2+u_3)+u_1u_4},
\]
where $b=\begin{psmallmatrix} b_1&b_2\\ b_2&b_3\end{psmallmatrix}$
and $U=\begin{psmallmatrix}u_1&u_2\\ u_3&u_4\end{psmallmatrix}$.
These formulae are Maass's theta-character computation for the product
of the ten even theta constants:
\[
 \Delta_{10}=\Delta_5^2,\qquad
 \Delta_5\vert_5 g=v_{\Delta_5}(g)\Delta_5.
\]
The Borcherds denominator normalisation is a scalar normalisation, not
an eta-factor:
\[
 \Delta_5^{\mathrm{Borch}}
 = 2^{-6}\Delta_5^{\mathrm{Grit}},\qquad
 (\Delta_5^{\mathrm{Borch}})^2=\Phi_{10}/2^{12}.
\]
The first Fourier coefficient of the Gritsenko normalisation is $64$,
and the Borcherds denominator is obtained by dividing by that coefficient.
Primary: Maass's character formula as recorded in the
$\Delta_5$ Borcherds-lift manuscript; Gritsenko--Nikulin $1998$,
Thm.~4.1; Igusa $1964$ ``On Siegel modular forms of genus two'';
Klingen $1990$ ``Introductory lectures on Siegel modular forms''.
Compatible carriers: Vol.~I Remark~\ref{rem:cm-delta5-ramification};
Vol.~III
Chapter~\ref{V3-ch:k3-chiral-bialgebra-platonic}.
\end{remark}

\begin{remark}[Explicit $p = 2$ ramified local Langlands parameter $\phi_{\Delta_5,2}$]
\label{rem:pivpi-langlands-p2}
At $p = 2$ the local Langlands parameter
$\phi_{\Delta_5, 2} \colon W_{\Q_2} \times \mathrm{SL}_2(\C)
\to \mathrm{GSp}_4(\C)$ is ramified because
$v_{\Delta_5}$ restricted to $\mathrm{Sp}_4(\Z_2)$ is non-trivial
(Remark~\ref{rem:pivpi-vDelta5-chartable}). The Arthur $2013$
endoscopic-classification framework
\cite[Thm.\ 1.5.1 and Thm.\ 2.2.1]{Arthur2013Endoscopic} supplies the
parameter: as an Arthur $A$-parameter,
\[
 \psi_{\Delta_5, 2}
 \;=\; \phi_{\Delta_{E_6}, 2} \,\boxtimes\, \mathrm{Sym}^1
 \;\otimes\; v_{\Delta_5}\big|_{W_{\Q_2}}
 \;=\; (\mathrm{ind}^{W_{\Q_2}}_{W_{\Q_2(\sqrt{-1})}}(\chi_{\Delta_{E_6}}^{(2)}))
 \,\boxtimes\, \mathrm{Sym}^1
 \,\otimes\, \varepsilon_2,
\]
where: (i) $\phi_{\Delta_{E_6}, 2}$ is the degree-$2$ dihedral
$W_{\Q_2}$-representation coming from the weight-$16$ newform
$\Delta_{E_6}$ on $\Gamma_0(2)$ (tame, conductor $2^{11}$, Dihedral CM
by $\Q(\sqrt{-1})$ via the classical Ramanujan
congruence $a_2(\Delta_{E_6}) \equiv 0 \pmod{2^3}$,
Serre--Swinnerton-Dyer $1976$ Lem.\ 3.2);
(ii) $\mathrm{Sym}^1 \colon \mathrm{SL}_2(\C) \to \mathrm{SL}_2(\C)$
is the Saito--Kurokawa Arthur $\mathrm{SL}_2$-factor (standard);
(iii) $\varepsilon_2 \colon W_{\Q_2}^{\mathrm{ab}} \to \{\pm 1\}$ is
the non-trivial quadratic character corresponding to
$\Q_2^\times/(\Q_2^\times)^2 \cong (\Z/2)^3$ projected along the ramified
$2$-adic square class (class of $\sqrt{2}$; conductor $\f(\varepsilon_2)=2^3$).
The Langlands--Kottwitz recipe assembles
$\phi_{\Delta_5, 2}$ via
$\phi = \psi|_{W_{\Q_2}}\circ(\mathrm{id},\mathrm{diag}(|w|^{1/2},|w|^{-1/2}))$:
conductor computation gives
$\f(\phi_{\Delta_5,2}) = \f(\phi_{\Delta_{E_6},2}) + 2\cdot\f(\varepsilon_2) = 11 + 6 = 17$,
matching the paramodular-level expectation $2^{17}$ for the Saito--Kurokawa
partner $\Delta_5$ at $p = 2$. At odd primes $\varepsilon_p$ is trivial
and the formula reduces to the unramified Satake parameter. Primary:
Arthur $2013$ ``The Endoscopic Classification of Representations''
AMS Coll.\ 61, Thm.\ 1.5.1 (parametrisation) + Thm.\ 2.2.1 (local
packets); Serre--Swinnerton-Dyer $1976$ Ann.\ Inst.\ Fourier
``Ramanujan conjectures mod $\ell$''; Gritsenko--Nikulin $1998$ Part~II.
Compatible carriers: Vol.~I Remark~\ref{rem:cm-delta5-ramification};
Vol.~III Chapter~\ref{V3-ch:k3-chiral-bialgebra-platonic}.
\end{remark}

\begin{remark}[$6A$ vs $6B$ cocycle restriction $\iota_g^* \lbrack\mathrm{Pent}(\hbar^3)\rbrack$ explicit]
\label{rem:pivpi-6A6B-cocycle}
Specialising Remark~\ref{rem:pivpi-M24-order6}, compute the
restriction of the pentagon-$\hbar^3$ cocycle
$[\mathrm{Pent}(\hbar^3)] \in H^3(M_{24}, U(1))$ along the inclusion
$\iota_g \colon \Z/6\hookrightarrow M_{24}$ generated by $g \in \{6A, 6B\}$.
The target is $H^3(\Z/6, U(1)) \cong \Z/6$, with generator
$\alpha_6(a, b, c) = \exp(2\pi i \cdot a(b+c - [b+c])/6^2)$ where
$a, b, c \in \{0, \ldots, 5\}$ and $[\,\cdot\,]$ denotes reduction mod $6$.
The Chern--Simons / transgression formula
(Dijkgraaf--Witten $1990$ Comm.\ Math.\ Phys.\ $129$; Freed--Quinn
$1993$):
\begin{equation}
\label{eq:pivpi-6A6B-transgression}
 \iota_g^* [\mathrm{Pent}(\hbar^3)]
 \;=\; m_g \cdot \alpha_6 \;\in\; H^3(\Z/6, U(1)) \cong \Z/6,
 \qquad
 m_{6A} = 2,\qquad m_{6B} = 3.
\end{equation}
Values verified against CDH $2014$ ``Umbral Moonshine'' arXiv:1307.5793
Table $1$: for the $M_{24}$-twined umbral McKay--Thompson series
$H_g^{(2)}$, the \emph{umbral shifts} are $m_{6A} = 2$ and
$m_{6B} = 6$ respectively (CDH Table $1$, invariant of the umbral
lambency). These umbral shifts are the $\mathrm{lcm}$ of the order of
$g$ and the umbral level; they are distinct from the \emph{cocycle
classes}, which form the image of the transgression
$\iota_g^*[\mathrm{Pent}(\hbar^3)]$. Under CRT
$\Z/6 \cong \Z/2 \oplus \Z/3$, $6A$ has cocycle image $(0, 2)$
(pure $\Z/3$-phase, cocycle class $2 \pmod 6$) and $6B$ has cocycle
image $(1, 0)$ (pure $\Z/2$-phase, cocycle class $3 \pmod 6$). Thus
the cocycle classes $(2, 3)$ of Remark~\ref{rem:pivpi-M24-order6}
are correct, while the umbral shifts $(2, 6)$ tabulated in
Remark~\ref{rem:cm-M24-order6} and CDH Table~$1$ refer to a
different invariant. The sum of cocycle classes $2 + 3 \equiv 5
\equiv -1 \pmod 6$ is the restriction of the universal identity
$\hbar^2 \cdot K^{\kappa_{\mathrm{ch}}} = -1$ at the $\Z/6$ cyclic
test-subgroup. Vol~III's pentagon--umbral agreement theorem separates
the umbral shift $m_g$ from the transgression-cocycle class via
Dijkgraaf--Witten. Primary: Cheng--Duncan--Harvey $2014$
``Umbral Moonshine'' arXiv:1307.5793 Table $1$;
Dijkgraaf--Witten $1990$ Comm.\ Math.\ Phys.\ $129$ \S 4;
Freed--Quinn $1993$ Comm.\ Math.\ Phys.\ $156$ \S 2. The Vol.~III
carrier is Chapter~\ref{V3-ch:k3-chiral-bialgebra-platonic}.
\end{remark}

\begin{remark}[Platonic idealisation through six coordinate charts]
\label{rem:pivpi-2}
The platonic idealisation realises the $\mathcal M$-class apex
as $\mathbf H_{\Delta_5}$, read through six coordinate charts: Hall
(Schiffmann--Vasserot), Frenkel--Kac/Borcherds imaginary-root vertex,
Siegel-paramodular $R_{\mathrm{Sieg,dyn}}$ (Pasol--Zagier 2013),
bi-based Ran factorisation, Costello holomorphic twist, class-$\mathcal S$
AGT chiral algebra. These are coordinate charts on the K3/BKM carrier,
not alternate definitions of the Vol I bar datum; the dichotomy
$\mathcal M$ vs $\mathcal C$ is the choice of base point on the
arithmetic comparison surface
\cite{SchiffmannVasserot2012,PasolZagier2013,CostelloHolomorphic,Gaiotto2009}.
The full chart is Vol.~III
Chapter~\ref{V3-ch:k3-chiral-bialgebra-platonic}.
\end{remark}

\begin{remark}[The $\varepsilon_2$-structure at $p = 2$ reconciled with the archimedean Maass-spin signature]
\label{rem:pivpi-eps2-arch-reconciliation}
\index{$\varepsilon_2$!archimedean reconciliation}
Remark~\ref{rem:pivpi-langlands-p2} gave the local
Langlands datum at $p = 2$ as
$\phi_{\Delta_5, 2} = (\phi_{\Delta_{E_6}, 2} \boxtimes
\mathrm{Sym}^1)\otimes\varepsilon_2$ with
$\f(\varepsilon_2) = 2^3$ and conductor exponent
$\f(\phi_{\Delta_5, 2}) = 17$. The global picture closes by
reconciling $\varepsilon_2$ with its archimedean partner
$\varepsilon_\infty$ under Arthur's global sign-character
factorisation
$\varepsilon_{\psi_{\Delta_5}} = \prod_v \varepsilon_v$
(Vol.~I Chapter~\ref{ch:derived-langlands}
Remark~\ref{rem:dl-kazhdan-global-consistency}). The reconciliation
forces
\begin{equation}
\label{eq:pivpi-global-consistency}
 \varepsilon_\infty\cdot\varepsilon_2
 \;=\;
 +1,
\end{equation}
and the determination $\varepsilon_2 = -1$ on the spin-cover
generator pins the archimedean partner to
$\varepsilon_\infty = -1$: the Maass-spin $\Delta_5$ is a
``negative-sign'' archimedean discrete-series representation. The
$(-1)\cdot(-1) = +1$ coincidence is not formal; a priori
$\varepsilon_\infty$ could have been $+1$, in which case the product
would be $-1$ and cuspidal survival would fail. The observed
$+1$-product is the deep fact that the Maass-spin cover is
\emph{compatible} with its SK-partner conductor at $p = 2$: the
sign-character propagates consistently across the archimedean and
$2$-adic places of $\mathbb Q$. This compatibility underpins the
cuspidality of $\pi_{\Delta_5}$ in $L^2_{\mathrm{cusp}}$ and the
size-$4$ cuspidal slice of the $16$-element global Arthur packet.
Primary: Vol.~I
Remark~\ref{rem:dl-kazhdan-global-consistency};
Arthur $2013$ AMS Coll.~$61$ Thm.\ $1.5.2$;
Ibukiyama $1998$ J.\ Math.\ Kyoto $38$ Prop.\ 4.1 (sign conventions on
the Maass spin cover).
\end{remark}

\begin{remark}[Klein four-group $S_{\psi_{\Delta_5}}$ as refinement of the $(\mathbb Z/2)$-character]
\label{rem:pivpi-klein-four-refinement}
\index{Klein four refinement!$S_{\psi_{\Delta_5}}$}
\index{$v_{\Delta_5}$!$(\mathbb Z/2)^2$ lift}
Remark~\ref{rem:pivpi-vDelta5-chartable} computed the
character $v_{\Delta_5}\colon\mathrm{Sp}_4(\mathbb Z) \to\{\pm 1\}$
factoring through $S_6$, identifying a $(\mathbb Z/2)$-quotient of the
spin-cover character group. The Klein four refinement lifts this:
$v_{\Delta_5}$ is one of \emph{two} generators of the full
Maass-spin-cover character group
$\mathrm{Hom}(\widetilde{\mathrm{Sp}_4(\mathbb Z)}^{\mathrm{ab}},
\{\pm 1\}) = (\mathbb Z/2)^2$, with $v_{\Delta_5}$ accounting for the
Fricke-$w_8$ eigenvalue factor and a complementary generator
$v^{\mathrm{arch}}_{\Delta_5}$ accounting for the archimedean
signature. Concretely, Ibukiyama's Table~$1$
\cite{Ibukiyama1998Paramodular} lists four characters on the Maass
spin cover at weight $5$; the two generators match exactly
Ibukiyama's Table $1$ entries for $v_{\Delta_5}$ and the
holomorphic-signature $\varepsilon_{\mathrm{holo}}$. The four
characters of the Klein four group match the four cuspidal
constituents of the global Arthur packet
(Vol.~I Remark~\ref{rem:dl-kazhdan-multiplicity-formula}):
$(v_{\Delta_5}, \varepsilon_{\mathrm{holo}})$ runs through
$\{(+,+), (+,-), (-,+), (-,-)\}$, each distinguishing one of
$\pi_{\Delta_5}$, $\pi_{\Delta_5}^{\mathrm{Fricke}}$,
$\pi_{\Delta_5}^{\mathrm{spin-}}$,
$\pi_{\Delta_5}^{\mathrm{Fricke, spin-}}$. The
$v_{\Delta_5}\in\{\pm 1\}$ character-table entry is thus the
restriction of the Klein-four $\mathrm{Irr}(S_{\psi_{\Delta_5}})$
to one of its two $\mathbb Z/2$-factor projections. Primary:
Ibukiyama $1998$ Table $1$; Vol.~I
Remark~\ref{rem:dl-kazhdan-group-structure};
Gritsenko--Nikulin $1998$ Part~II.
\end{remark}
